\vspace{.5cm}
\begin{table}[H]
\caption{Comparisons of Alternative Policies}
\label{tab:comparison_1}
\centering
\begin{threeparttable}
    \begin{tabular*}{\textwidth}{l @{\extracolsep{\fill}} cc}
        \toprule
        & Difference in welfare gains & p-value\\
        \midrule
%        \addlinespace
       Dynamic targeting ($\pi^{\ast}$) vs. 100\% $(T,U)$ & $1365.7$ & $0.000$ \\
        & $(309.1)$ & \\
        \addlinespace
        Dynamic targeting ($\pi^{\ast}$) vs. 100\% $(U,T)$ & $1546.5$ & $0.000$ \\
        & $(328.7)$ & \\
        \addlinespace
        Dynamic targeting ($\pi^{\ast}$) vs. 100\% $(T,T)$ & $1397.7$ & $0.000$ \\
        & $(319.8)$ & \\
        \addlinespace
       Dynamic targeting ($\pi^{\ast}$) vs. Static targeting I ($\pi_{S(I)}^{\ast}$) & $913.8$ & $0.000$ \\
        & $(269.2)$ & \\
        \addlinespace
        Dynamic targeting ($\pi^{\ast}$) vs. Static targeting II ($\pi_{S(\two)}^{\ast}$)& $839.1$ & $0.002$ \\
        & $(287.8)$ & \\
        \bottomrule
    \end{tabular*}
\begin{tablenotes}[flushleft]
\footnotesize\item Notes: This table compares welfare gains from each policy. For each row, the column ``Difference in Welfare Gains'' shows the estimated welfare gain of the policy on the left-hand side ($W_L$) relative to the policy on the right-hand side ($W_R$) in JPY per household, with its standard error in parenthesis. The column ``p-value'' gives the p-value for the null hypothesis: $H_0: W_L \leq W_R$. The monetary unit is given as 1 ¢ = 1 JPY, the exchange rate in the summer of 2020.
\end{tablenotes}
\end{threeparttable}
\end{table}
\vspace{.5cm}